\documentclass[10pt,a4paper]{amsart}
\usepackage[margin=19mm]{geometry}
\usepackage{amsmath,amssymb,mathtools}
\usepackage[scr=rsfs,cal=euler]{mathalfa}
\usepackage{xfrac}
\usepackage[unicode,hidelinks,bookmarks=false]{hyperref}
\newcommand{\ie}{i.e.\ }
\newcommand{\cf}{cf.\ }
\newcommand{\resp}{respectively\ }
\newcommand{\Subsection}{\subsection}
\providecommand{\nobreakdash}{\nobreak\hbox{-}}
\setlength{\emergencystretch}{2em}
\multlinegap0pt
\usepackage{amssymb}
\usepackage{multirow}
\usepackage[all]{xy}
\usepackage{enumerate}
\usepackage{mathtools}
\usepackage{mdframed}
\usepackage{booktabs}
\usepackage{caption}
\newtheorem{definition}{Definition}
\newtheorem{theorem}{Theorem}
\newtheorem{lemma}{Lemma}
\newtheorem{proposition}{Proposition}
\title{Equivariant Koszul Cohomology of Canonical Curves}
\author{Kostas Karagiannis, Aristides Kontogeorgis, Konstantia Manousou Sotiropoulou}
\date{Source: arXiv:2512.19894v2 (CC BY 4.0)}
\begin{document}
\maketitle

\noindent\emph{Source and licence.} Equivariant Koszul Cohomology of Canonical Curves. Version arXiv:2512.19894v2 (CC BY 4.0), distributed by arXiv under the Creative Commons Attribution 4.0 International licence, http://creativecommons.org/licenses/by/4.0/, as recorded in the arXiv licence metadata for this exact version and shown on the abstract page https://arxiv.org/abs/2512.19894v2. Full licence notice: \texttt{licenses/arxiv-2512.19894-CC-BY-4.0.txt}.\par\smallskip

\noindent\textit{Editorial scope.} Counted unit: the article's theorem th:mainwedge (source0.tex lines 1012--1026). The article's complete original proof of it is delivered with it (source0.tex lines 1027--1183, one \texttt{proof} environment of 157 lines). No mathematics is written, repaired or re-derived: the statement and the proof are the article's own lines, in the article's own notation, and no retained line is substituted or re-flowed.  Retained uncounted as the source-local context the proof needs: lines 662--679; lines 701--719; lines 775--781; lines 820--826; lines 853--862; lines 915--929; lines 978--982; lines 984--995.  Nothing was omitted from any retained span, so the package carries no omission record.  No retained line contains a citation, so the wrapper carries no bibliography and no bibliography entry.  No retained line contains a figure environment, an image-inclusion command or drawing code, so no image is delivered and the package declares no asset dependency.  Editorial formatting only, in addition: an AMS \texttt{amsart} preamble that restates the article's own declarations and macros used by the retained lines, namely the environments \texttt{definition}, \texttt{theorem}, \texttt{lemma}, \texttt{proposition} on separate counters, together with the standard packages those lines need.  The retained environments print in this document as Definition 1, Theorem 1, Lemma 1, Proposition 1 in order of appearance, whereas the article prints the same items with the numbering of its own full document.  The complete native source of the article as distributed in the arXiv e-print for this exact version is bundled unchanged as \texttt{sources/191519/source0.tex}.
\begin{definition}\label{def:localmult}
Let $Q\in \mathfrak{b}$, let $P\in\pi^{-1}(Q)\subseteq\mathfrak{r}$ and let $\mathcal{E}\in{\rm QCoh}_G(X)$ be a vector bundle. 
\begin{enumerate}
\item The completed fiber of $\mathcal{E}$ over $P$ is the $kG_P$-module
\[
V_{P}(\mathcal{E})=(\mathcal{E}\otimes_{C_Q}\widehat{\mathcal{O}}_{X,P})\otimes_{\widehat{\mathcal{O}}_{X,P}} k,
\text{ where }
C_Q=\mathcal{O}_X\otimes_{\mathcal{O}_Y}\widehat{\mathcal{O}}_{Y,Q}.
\]
\item The local multiplicities of an irreducible $kG$-module $I$ with respect to $\mathcal{E}$ are the integers
\[
n_{d,Q,I}(\mathcal{E})=\langle \chi_P^d, V_P(\mathcal{E})\otimes {\rm Res}^G_{G_P}(I)\rangle,
\text{ where }
0\leq d\leq e_Q-1.
\]
In the special case that $\mathcal{E}=\mathcal{O}_X$, we write $n_{d,Q,I}$ for $n_{d,Q,I}(\mathcal{O}_X)=\langle \chi_P^d, {\rm Res}^G_{G_P}(I)\rangle$.
\end{enumerate}
\end{definition}

\begin{theorem} \label{th:EL}
If $\mathcal{E}\in{\rm QCoh}_G(X)$ is a vector bundle on $X$, then
\[
\chi_G(\mathcal{E})
=
  \left( 
\frac{1}{|G|} \deg(\mathcal{E})+(1-g_Y) \mathrm{rk}(\mathcal{E})
  \right)[kG]
  - 
  \sum_{I \in \mathrm{Irr}(G)} 
  \sum_{Q\in Y}
  \sum_{d=1}^{e_Q}
\left(
1 -\frac{d}{e_Q}
\right)
n_{d,Q,I}(\mathcal{E}) [I],
\]
where $Y=X/G$, $g_Y$ is the genus of $Y$, ${\rm Irr}(G)$ is the set of irreducible representations of $G$, $e_Q$ is the ramification index of $P\in\pi^{-1}(Q)$ and the $n_{d,Q,I}(\mathcal{E})$ are as in Definition \ref{def:localmult}. 
\end{theorem}

\begin{lemma}\label{lem:nomega}
For all $Q\in Y,\;1\leq d\leq e_Q,\; I\in{\rm Irr}(G)$ and $m\in\mathbb{N}$,
\[
n_{d,Q,I}\big(\Omega_X^{\otimes m}\big)=n_{e_Q\left\{\frac{m-d}{e_Q}\right\},Q,I^\vee},
\]
where $\{ x\}=x-\lfloor x \rfloor$ denotes the fractional part of $x\in\mathbb{Z}$ and $I^\vee$ is the dual of $I$.
\end{lemma}

\begin{proposition}\label{prop:vanish}
$
H^1\big(\bigwedge^p\mathcal{M}_{\Omega_X}\otimes \Omega_X^{\otimes 2}\big)=0
$ if $p\neq g-1$, and $
H^1\big(\bigwedge^{g-1}\mathcal{M}_{\Omega_X}\otimes \Omega_X^{\otimes 2}\big)=k
$.
\end{proposition}

\begin{proposition}\label{prop:rkdeg}
\[
{\rm rk}\big(\bigwedge^p\mathcal{M}_{\Omega_X}\otimes \Omega_X^{\otimes 2}\big)
=
\binom{g-1}{p}
\text{ and }
\deg( \bigwedge^p M_{\Omega_X} \otimes \Omega_X^q)=
2(2g-2-p)\binom{g-1}{p}.
\]
\end{proposition}

\begin{proposition}\label{prop:binomial3}
If $Q\in Y,\;0\leq d\leq e_Q-1,\;I\in{\rm Irr}(G)$ and $p\in\mathbb{Z}_{\geq 0}$, then
\[
n_{d,Q,I}\big(\bigwedge^p\mathcal{M}_{\Omega_X}\otimes \Omega_X^{\otimes 2}\big)
=\sum_{i=0}^p
\sum_{k=0}^{e_Q-1}
(-1)^i
\langle
\chi_P^{d-k},
\bigwedge^{p-i}H^0(\Omega_X)
\rangle
\;
n_{k,Q,I}(\Omega_X^{\otimes (i+2)}).
\]
\end{proposition}

\begin{equation}\label{eq:ChW}
[H^0\big(\Omega_X\big)]=\sum_{I\in{\rm Irr}(G)} m_I[I],
\text{ where }
m_I=(g_Y-1)\dim I+\sum_{Q\in Y}\sum_{d=0}^{e_Q-1}\left\langle \frac{-d}{e_Q}\right\rangle n_{d,Q,I}+\delta_{I,I_{\rm triv}}
\end{equation}

\begin{lemma}\label{lem:localwedge}
For $Q\in Y,\;0\leq d\leq e_Q-1,\;r\in\mathbb{Z}_{\geq 0}$ and $\{m_I:\;I\in {\rm Irr}(G)\}$ as in eq. (\ref{eq:ChW}),
\[
\langle 
\chi_P^{d}, \bigwedge^{r}H^0(\Omega_X)
\rangle
=
\sum_{\substack{ 0\leq k_j\leq \sum_{I}m_I n_{j,Q,I}  \\ \sum k_j=r \\ \sum jk_j\equiv d\mod e_Q}}
\prod_{j=0}^{e_Q-1}
\left(\!\begin{array}{c} \displaystyle\sum_{I\in{\rm Irr}(G)}m_I\cdot n_{j,Q,I} \\ k_j \end{array}\!\right).
\]
\end{lemma}

\begin{theorem}\label{th:mainwedge}
Let $g_X$ be the genus of $X$ and $g_Y$ the genus of $Y=X/G$. For each $Q\in Y$, let $e_Q$ be the ramification index of $P\in\pi^{-1}(Q)$, let $G_P$ be its decomposition group, $\chi_P$ be the fundamental character of $\mathfrak{m}_P/\mathfrak{m}_P^2$ and let $n_{k,Q,I}=\langle \chi_P^k,I\rangle$, for $I\in{\rm Irr}(G), \;0\leq k\leq e_Q-1$. Then
\begin{multline*}
[H^0\big(\bigwedge^p\mathcal{M}_{\Omega_X}\otimes\Omega_X^{\otimes 2}\big)]
=
\binom{g_X-1}{p} \left( \frac{2(g_X-1-p)}{|G|}
+
g_Y-1\right)
[kG]\\
  +
\sum_{I \in \operatorname{Irr}(G)} \sum_{Q \in Y} \sum_{i=0}^p (-1)^i \sum_{d=0}^{e_Q-1} 
\langle \chi_P^{d}, \bigwedge^{p-i}H^0(\Omega_X) \rangle 
\sum_{k=0}^{e_Q-1} \left\{ \frac{d+i+1-k}{e_Q} \right\} n_{k,Q,I} [I].
\end{multline*}
\end{theorem}

\begin{proof}
Combining Theorem (\ref{th:EL}) and Proposition \ref{prop:vanish}, gives

\begin{align*}
[H^0\big(\bigwedge^p\mathcal{M}_{\Omega_X}\otimes\Omega_X^{\otimes 2}\big)]
=&
  \left( 
\frac{1}{|G|} \deg\big(\bigwedge^p\mathcal{M}_{\Omega_X}\otimes\Omega_X^{\otimes 2}\big)+(1-g_Y) \mathrm{rk}\big(\bigwedge^p\mathcal{M}_{\Omega_X}\otimes\Omega_X^{\otimes 2}\big)
  \right)[kG]\\
  &- 
  \sum_{I \in \mathrm{Irr}(G)} 
  \sum_{Q\in Y}
  \sum_{d=1}^{e_Q}
\left(
1 -\frac{d}{e_Q}
\right)
n_{d,Q,I}\big(\bigwedge^p\mathcal{M}_{\Omega_X}\otimes\Omega_X^{\otimes 2}\big) \cdot [I].
\end{align*}
Substituting the degree and the rank from Proposition \ref{prop:rkdeg} into the coefficient of $[kG]$ gives
\begin{equation}\label{eq:rkdeg}
\binom{g_X-1}{p} \left( 
\frac{2(2g_X-2-p)}{|G|} +(1-g_Y)
  \right)[kG].
\end{equation}
Regarding the triple sum, Proposition \ref{prop:binomial3} and Lemma \ref{lem:nomega} imply that
\[
\left(1-\frac{d}{e_Q}\right)n_{d,Q,I}\big(\bigwedge^p\mathcal{M}_{\Omega_X}\otimes \Omega_X^{\otimes 2}\big)
=
\sum_{i=0}^{p}\sum_{k=0}^{e_Q-1}
(-1)^i
\left(1-\frac{d}{e_Q}\right)
\;
\langle
\chi_P^{d-k},
\bigwedge^{p-i}H^0(\Omega_X)
\rangle
\;
n_{e_Q\left\{ \frac{i+2-k}{e_Q}\right\},Q,I^\vee},
\]
and observing that $k'=e_Q\left\{ \frac{i+2-k}{e_Q}\right\}$ if and only if $k=e_Q\left\{ \frac{i+2-k'}{e_Q}\right\}$ allows us to rewrite the right hand side
\[
\sum_{i=0}^{p}\sum_{k=0}^{e_Q-1}
(-1)^i
\left(1-\frac{d}{e_Q}\right)
\;
\langle
\chi_P^{d-e_Q\left\{ \frac{i+2-k}{e_Q}\right\}},
\bigwedge^{p-i}H^0(\Omega_X)
\rangle
\;
n_{k,Q,I^\vee}.
\]
Finally, $d-e_Q\left\{ \frac{i+2-k}{e_Q}\right\}\equiv e_Q\left\{ \frac{d-i-2+k}{e_Q}\right\}{\rm mod }\;e_Q$. Setting  $d'
=
e_Q
\left\{
\frac{d-i-2+k}{e_Q}
\right\}
$ forces
$
d=
1+
e_Q
\left\{
\frac{d'+i+1-k}{e_Q}
\right\}$,
to ensure that $1\leq d \leq e_Q$ with $0\leq d'\leq e_Q-1$.
Thus, we arrive at
\begin{multline*}
\sum_{d=1}^{e_Q}\left(1-\frac{d}{e_Q}\right)n_{d,Q,I}\big(\bigwedge^p\mathcal{M}_{\Omega_X}\otimes \Omega_X^{\otimes 2}\big)
=
\\
\sum_{d=0}^{e_Q-1}\sum_{i=0}^{p}\sum_{k=0}^{e_Q-1}
(-1)^i
\left(
1-\frac{1}{e_Q}-
\left\{
\frac{d+i+1-k}{e_Q}
\right\}
\right)
\;
\langle
\chi_P^{d},
\bigwedge^{p-i}H^0(\Omega_X)
\rangle
\;
n_{k,Q,I^\vee}.
\end{multline*}
The term $(1-\frac{1}{e_Q})$ is independent of the summation indices, and so we write the RHS as
\begin{multline}\label{eq:crazysum}
\left(1-\frac{1}{e_Q}\right)
\sum_{i=0}^{p}
(-1)^i
\sum_{d=0}^{e_Q-1}
\langle
\chi_P^{d},
\bigwedge^{p-i}H^0(\Omega_X)
\rangle
\sum_{k=0}^{e_Q-1}
n_{k,Q,I^\vee}
\\
-
\sum_{d=0}^{e_Q-1}\sum_{i=0}^{p}\sum_{k=0}^{e_Q-1}
(-1)^i
\left\{
\frac{d+i+1-k}{e_Q}
\right\}
\;
\langle
\chi_P^{d},
\bigwedge^{p-i}H^0(\Omega_X)
\rangle
\;
n_{k,Q,I^\vee}.
\end{multline}
Next, we observe that $\sum_k n_{k,Q,I^\vee}=\dim I^\vee$ and
\[
\sum_{i=0}^{p}
(-1)^i
\sum_{d=0}^{e_Q-1}
\langle
\chi_P^{d},
\bigwedge^{p-i}H^0(\Omega_X)
\rangle
=
\sum_{i=0}^{p}
(-1)^i \binom{g_X}{p-i}
=
\binom{g_X-1}{p}
\]
and so
\[
\sum_{Q\in Y}\sum_{I\in{\rm Irr}(G)} \left(1-\frac{1}{e_Q}\right)
\sum_{i=0}^{p}
(-1)^i
\sum_{d=0}^{e_Q-1}
\langle
\chi_P^{d},
\bigwedge^{p-i}H^0(\Omega_X)
\rangle
\sum_{k=0}^{e_Q-1}
n_{k,Q,I^\vee}
[I]
=
\sum_{Q\in Y}\left(1-\frac{1}{e_Q}\right)
\binom{g_X-1}{p}
[kG].
\]
Using the Riemann-Hurwitz formula, the RHS becomes
\[
\binom{g_X-1}{p}\left(\frac{2(g_X-1)}{|G|}-2(g_Y-1)\right)[kG]
\]
and combining this with eq. (\ref{eq:rkdeg}) gives the final coefficient of $[kG]$. The remaining term is obtained from the second line of eq. (\ref{eq:crazysum}). The reader should keep in mind that an explicit formula for the integers $\langle
\chi_P^{d},
\bigwedge^{p-i}H^0(\Omega_X)
\rangle$ has been given in Lemma \ref{lem:localwedge}.
\end{proof}

\end{document}
